\documentclass[titlepage, 10pt]{ltjsarticle}
\usepackage[utf8]{inputenc}
\usepackage{graphicx, color, xcolor}
\usepackage{ascmac}
\usepackage{float}
\usepackage{enumitem}
\usepackage{tabularray}
\usepackage{listings, jvlisting}
\usepackage{amssymb}
\usepackage{wrapfig}
\usepackage{diagbox}
\usepackage{siunitx}
\usepackage{enshuc}

\renewcommand{\subsubsection}{・}

\title{情報科学演習C \\ 課題3} 
\author{植村 向陽}
\studentid{09B24008}
\affiliation{基礎工学部情報科学科ソフトウェア科学コース}
\date{\today} 
\begin{document}
\maketitle 
\lstset{
    language=C,
    basicstyle=\ttfamily\small,
    numbers=left,
    numberstyle=\tiny,
    stepnumber=1,
    numbersep=5pt,
    frame=single,
    breaklines=true,
}

\section{chatserver}
チャットサーバーのサーバープログラムとしてchatserver.cを作成した
\subsection{データ構造}
サーバー側に各クライアントのソケットとユーザーネームを格納するデータ構造として以下の構造体を定義した．
\begin{lstlisting}
typedef struct {
    int socket;
    char name[MAX_USERNAME_LENGTH];
    ClientStatus status;
} ClientInfo;

typedef struct _ClientList {
    ClientInfo clients[MAXCLIENT];
    int client_count;
    int max_client;
} ClientList;

\end{lstlisting}
\section{実行結果}
\end{document}